\documentclass[journal]{IEEEtran}

\usepackage{amsmath,amssymb}
\usepackage{booktabs}
\usepackage{cite}
\usepackage{graphicx}
\usepackage{multirow}
\usepackage{url}
\usepackage{xcolor}

\newcommand{\method}{VANET-FL-Former}
\newcommand{\todoauthor}{Anonymous Author(s)}
\newcommand{\revisionnote}[1]{\textcolor{red!70!black}{\textbf{Submission requirement:} #1}}

\begin{document}

\title{Auditing Federated Transformer Intrusion Detection in VANETs:\\
27-Class Performance Under Noisy Clients}

\author{\todoauthor%
\thanks{Author names, affiliations, funding, and the corresponding-author address are withheld for double-blind review.}}

\markboth{Manuscript under preparation}%
{Anonymous Author(s): Auditing Federated Transformer Intrusion Detection in VANETs}

\maketitle

\begin{abstract}
Federated learning (FL) is attractive for vehicular intrusion detection because safety messages can remain at the vehicle or roadside unit while model updates are shared. Its evaluation is nevertheless easy to overstate: network attacks in the intrusion-detection corpus are distinct from malicious FL clients, a nominally robust aggregator may perform no trimming at small client counts, and a token transformer is not necessarily a temporal model across messages. This paper presents an evidence-driven audit and preliminary evaluation of \method, a compact transformer-based, 27-class FL intrusion detector implemented in Flower. The balanced VANET-IDS26 corpus contains 27,000 records (1,000 per class), including benign traffic and 26 attack classes. Records are distributed across four clients with a label-Dirichlet partition; each client trains a 1.50-million-parameter encoder locally for two epochs in each of ten communication rounds. In the evaluated adversarial condition, one of four clients adds Gaussian noise with standard deviation 0.05 to every trained parameter before upload. The final global model obtains 77.20\% accuracy, 0.77 macro F1, and 0.5848 cross-entropy on a 5,400-record federated holdout. Recall is perfect for false-emergency-vehicle, flooding-DDoS, impossible-kinematics, and Sybil attacks, but falls to 0.40 for heading manipulation and 0.44 for position offset. The audit also finds that $\lfloor\beta K\rfloor=0$ for $K=4$ clients and trim proportion $\beta=0.1$, making the configured trimmed mean equivalent to an unweighted coordinate-wise mean. The 64-token representation contains only the first 62 data tokens of one serialized message, cutting off most claimed-state and all event-type fields; evaluation uses per-client random holdouts rather than the canonical test split. We therefore report the measured result as a reproducible preliminary benchmark, not as evidence of Byzantine robustness, temporal learning, or external generalization. The paper concludes with a controlled protocol required to establish these claims.
\end{abstract}

\begin{IEEEkeywords}
Byzantine robustness, federated learning, intrusion detection, Internet of Vehicles, model poisoning, transformer, vehicular ad hoc network.
\end{IEEEkeywords}

\section{Introduction}
\IEEEPARstart{V}{ehicular} ad hoc networks (VANETs) exchange safety-critical state among vehicles and infrastructure. Falsified positions, identities, timestamps, hazards, and kinematics can therefore affect both cyber security and physical safety. Machine-learning intrusion detection systems (IDSs) can learn these heterogeneous attack signatures, but central collection of mobility traces creates privacy, governance, and scaling concerns. Federated learning (FL) addresses the data-locality aspect by aggregating locally trained models rather than raw records \cite{mcmahan2017communication,kairouz2021advances}. Flower provides a practical framework for implementing and evaluating this process across heterogeneous clients \cite{beutel2020flower}.

FL does not by itself secure the learning process. A compromised participant can upload arbitrary parameters, and robust aggregation has assumptions that are easy to violate in small or non-identically distributed deployments \cite{yin2018byzantine,fang2020local}. This distinction is especially important for an IDS: an \emph{attack sample} is a supervised example that the detector should learn, whereas a \emph{malicious FL client} attacks collaborative training. Setting a client option to \texttt{none} removes client-side poisoning; it does not remove malicious VANET traffic from the corpus.

Recent VANET IDS work has combined FL with BERT-style sequence models \cite{ahsan2024privacy}, hybrid packet/physics models \cite{chen2024fast}, and blockchain-supported coordination \cite{mansouri2025distributed}. Transformers are appealing because self-attention can represent interactions among serialized message fields \cite{vaswani2017attention,devlin2019bert}. However, calling such a model temporal is justified only when its input contains multiple time-ordered messages or explicit temporal state. A transformer over tokens from one message captures intra-record order, not longitudinal behavior.

This work reports and audits a complete four-client experiment over VANET-IDS26. Rather than treating one headline score as sufficient, we connect the metric to the exact evaluation set, poisoning operation, aggregation implementation, and model input. The contributions are:
\begin{itemize}
    \item A reproducible 27-class FL pipeline that serializes vehicular telemetry, trains a compact transformer on four non-IID clients, and reports class-resolved performance.
    \item An adversarial-client experiment in which one client perturbs its locally trained parameters with Gaussian noise. The resulting global model reaches 77.20\% federated-holdout accuracy and 0.77 macro F1 after ten rounds.
    \item An implementation-level audit showing that the chosen four-client, $\beta=0.1$ trimmed-mean setting trims no update, and hence cannot support a claim of trimmed-mean robustness.
    \item A validity audit identifying severe field truncation, random rather than group/temporal holdouts, and the absence of cross-message temporal modeling. We translate these findings into a controlled evaluation protocol for a full journal study.
\end{itemize}

The purpose of these disclosures is not to discount the observed classification performance. It is to delimit exactly what that performance demonstrates and to make the next experiment falsifiable.

\section{Background and Related Work}

\subsection{Federated Optimization and Non-IID Data}
In synchronous cross-silo FL, server round $t$ broadcasts parameters $\mathbf{w}^{t}$ to a set of clients. Client $k$ minimizes its empirical objective $F_k$ for $E$ local epochs and returns $\mathbf{w}^{t+1}_k$. Federated Averaging combines local parameters in proportion to local sample counts \cite{mcmahan2017communication}. Vehicular clients naturally differ in geography, traffic density, and observed attacks, making local data non-IID. Dirichlet label partitioning is commonly used to create controlled label skew; lower concentration generally produces stronger heterogeneity \cite{li2022federated}.

The present implementation uses a label-Dirichlet partition and full participation. Unlike sample-weighted FedAvg, Flower's evaluated trimmed-mean routine computes a coordinate-wise mean over retained client values without weighting by local sample count. This matters because client shard sizes range from 5,885 to 7,515 records.

\subsection{Poisoning and Robust Aggregation}
Byzantine clients can send arbitrary updates. Coordinate-wise median and trimmed mean are classical defenses with statistical guarantees under explicit bounds on the adversarial fraction and distributional assumptions \cite{yin2018byzantine}. Krum instead selects an update based on distances to neighboring updates \cite{blanchard2017machine}. Later attacks showed that robust rules can still be manipulated when an adversary adapts to the aggregation mechanism \cite{fang2020local}; backdoor attacks further demonstrate that strong aggregate accuracy can coexist with targeted failures \cite{bagdasaryan2020backdoor}.

For $K$ scalar client values at a model coordinate, a symmetric trimmed mean with proportion $\beta$ removes
\begin{equation}
    b=\left\lfloor \beta K \right\rfloor
    \label{eq:trim-count}
\end{equation}
values from each tail and averages the remaining $K-2b$ values. Robustness therefore depends on the integer $b$, not only the displayed hyperparameter $\beta$.

\subsection{Learning-Based VANET Intrusion Detection}
VeReMi established a reproducible basis for comparing VANET misbehavior detectors \cite{vanderheijden2018veremi}. FL-BERT subsequently combined BERT-style classification with federated training for software-defined VANET intrusion detection \cite{ahsan2024privacy}. VAN-FED-IDS used recurrent packet analysis, physics-based features, and evidence fusion, and evaluated Flower-based deployments \cite{chen2024fast}. A distributed FL/blockchain framework similarly explored collaborative VANET detection and model integrity \cite{mansouri2025distributed}. These systems motivate privacy-aware collaborative IDSs, while differences in datasets, attack taxonomies, split units, and class counts make direct comparison of headline accuracy inappropriate.

Table~\ref{tab:related} positions the present work by evaluation scope rather than by incomparable accuracy.

\begin{table*}[t]
\centering
\caption{Scope of representative learning-based vehicular IDS studies. A check mark means the dimension is explicitly evaluated, not that all threats are solved.}
\label{tab:related}
\begin{tabular}{p{3.0cm}p{2.6cm}p{2.8cm}cccc}
\toprule
Study & Data/domain & Detection model & FL & Multi-class & Non-IID & Poisoned client \\
\midrule
van der Heijden \emph{et al.} \cite{vanderheijden2018veremi} & VeReMi/V2X & Misbehavior benchmark & -- & $\checkmark$ & -- & -- \\
Ahsan \emph{et al.} \cite{ahsan2024privacy} & VeReMi/VANET & FL-BERT & $\checkmark$ & $\checkmark$ & $\checkmark$ & -- \\
Chen \emph{et al.} \cite{chen2024fast} & F2MD/VANET & Bi-LSTM, LightGBM, fusion & $\checkmark$ & $\checkmark$ & $\checkmark$ & -- \\
Mansouri \emph{et al.} \cite{mansouri2025distributed} & VeReMi/VANET & FL models and blockchain & $\checkmark$ & $\checkmark$ & -- & -- \\
This study & VANET-IDS26/V2X & Compact transformer & $\checkmark$ & 27 classes & $\checkmark$ & Gaussian parameter noise \\
\bottomrule
\end{tabular}
\end{table*}

\section{System and Threat Model}

\subsection{Deployment Model}
We model four FL clients coordinated by one honest server. A client represents a logical vehicular data silo, such as a roadside unit or a simulation partition. All four clients participate in every round. Raw CSV records remain client-side during training; model parameters and scalar metrics are sent to the server through Flower. This is a data-locality architecture, not a complete privacy guarantee: parameters may leak information unless combined with secure aggregation or formal privacy mechanisms \cite{kairouz2021advances}.

\subsection{Adversary}
The server and three clients are honest. After the fourth client completes local AdamW training, it perturbs every model parameter:
\begin{equation}
    \widetilde{\mathbf{w}}_4^{t+1}
      = \mathbf{w}_4^{t+1} + \boldsymbol{\epsilon}^{t},
    \qquad
    \boldsymbol{\epsilon}^{t}\sim\mathcal{N}(\mathbf{0},\sigma^2\mathbf{I}),
    \label{eq:noise}
\end{equation}
where $\sigma=0.05$. The attacker does not adapt to other clients, the current global model, or a target class. This is a coarse availability-oriented model-poisoning stressor, not a stealthy backdoor or an optimal local-model poisoning attack.

\subsection{Nominal and Effective Aggregation}
The server is configured with Flower FedTrimmedAvg, $K=4$, and $\beta=0.1$. Equation~\eqref{eq:trim-count} gives $b=\lfloor0.4\rfloor=0$. Thus no high or low coordinate is removed and the effective update is
\begin{equation}
    \mathbf{w}^{t+1}=\frac{1}{4}\sum_{k=1}^{4}\widetilde{\mathbf{w}}_k^{t+1}.
    \label{eq:effective-mean}
\end{equation}
The result in this paper consequently measures tolerance of an unweighted mean to the specified noise at 25\% adversarial participation. It does \emph{not} validate trimmed-mean Byzantine resilience. With four clients, trimming one value from each tail requires $\beta\geq0.25$; that would leave only two values and should be assessed against coordinate-wise median and sample-weighted FedAvg.

\section{Data and Representation}

\subsection{VANET-IDS26 Corpus}
The source corpus \cite{vanetids26dataset} contains 27,000 records balanced exactly across 27 labels: benign and the 26 attack classes in Table~\ref{tab:classes}. The canonical manifest defines stratified training, validation, and test splits of 21,600, 2,700, and 2,700 records, respectively. Each class contributes 800/100/100 records. The large on-disk master additionally contains a deterministic payload field used as a storage artifact; this field is not tokenized by the model and must not be interpreted as additional behavioral evidence.

The evaluated client manifest was generated from all 27,000 records with seed 42 and Dirichlet concentration $\alpha=2.0$. The resulting shard sizes are 7,515, 6,715, 6,885, and 5,885. Every shard includes all 27 labels but in different proportions. Within each client, a deterministic random 80/20 split forms local training and holdout sets. Flooring the training boundary produces a combined 5,400-record holdout for the noisy-client run.

\subsection{Structured-Text Encoding}
Each record is serialized into key--value text, for example:
\begin{quote}\small
\texttt{time=125.50 ; claimed\_time=95.50 ; physical\_sender\_id=17 ; \ldots ; malicious\_delay\_ms=30000 ; event\_type=timestamp\_shift}
\end{quote}
Regex tokenization lowercases alphanumeric identifiers and punctuation. A shared vocabulary is built before client partition training; the preserved run contains 5,292 tokens. Sequences are truncated or padded to 64 tokens and surrounded by special classification and separator tokens.

\subsection{Input-Visibility and Temporal Audit}
The representation exposes two validity risks. First, sequence truncation is determined by serialization order rather than feature semantics. The 64-token sequence reserves two special tokens, leaving only the first 62 data tokens. In audited records, this span covers time, claimed time, sender identifiers, sequence number, true position/speed/acceleration/heading/lane, and only the beginning of claimed $x$. It excludes claimed $y$, claimed speed, acceleration, heading and lane, malicious delay, and all event/object fields. Consequently, several attack-defining discrepancies are partly or entirely invisible to the model.

The raw serialization does contain an exact class-name proxy through \texttt{event\_type} for nine classes and 9,000 records (33.33\%). However, these tokens occur at positions 98--102 (or later) and none survives the evaluated 62-data-token cutoff. There is therefore no direct event-name leakage in this run. The raw field remains a latent leakage risk if sequence length or field order changes, so a submission implementation should remove post-event metadata rather than rely on truncation for exclusion. Early sender and sequence identifiers can still support simulation-specific memorization under a random record split; sender- and run-disjoint evaluation is required.

Second, each sample is one serialized message. Positional embeddings capture token position within that message; there is no sliding window, recurrent state, sender history, or cross-message attention. We consequently describe the current model as a \emph{structured-message transformer}, not as evidence of temporal sequence learning.

\section{Model and Federated Training}

\subsection{Transformer Classifier}
Let $\mathbf{x}\in\mathbb{N}^{L}$ be a token sequence and $\mathbf{m}$ its attention mask. Token and learned position embeddings are summed after prepending a learned classification vector. Four pre-normalized transformer encoder blocks apply eight-head self-attention and a feed-forward network of width $4d$, where $d=128$. The final classification representation passes through layer normalization and a two-layer GELU head with dropout 0.15. Cross-entropy over 27 classes is minimized. The architecture contains 1,499,163 trainable parameters (approximately 5.72 MiB in FP32), much smaller than a pretrained BERT base model but trained from scratch.

\subsection{Training Protocol}
All clients initialize from the same server model. Each round performs two local epochs with batch size 16, AdamW learning rate $3\times10^{-4}$, weight decay $10^{-2}$, and cross-entropy loss. All clients then upload parameters; client 3 applies Eq.~\eqref{eq:noise} first. Training uses ten synchronous rounds, full client participation, and seed 42. Server evaluation reconstructs the global model and evaluates the merged local holdouts after every round.

Table~\ref{tab:configuration} records the complete evaluated configuration. The experiment ran with Python 3.12.11, PyTorch 2.12.0, Flower 1.32.1, scikit-learn 1.8.0, and an NVIDIA L40S GPU. All clients and the server execute on one host, so this is a logical federation; network latency, disconnection, and hardware heterogeneity are outside scope.

\begin{table}[t]
\centering
\caption{Evaluated configuration.}
\label{tab:configuration}
\begin{tabular}{ll}
\toprule
Component & Value \\
\midrule
Classes / records & 27 / 27,000 \\
Clients / participation & 4 / 100\% per round \\
Dirichlet concentration & $\alpha=2.0$ \\
Client records & 7,515; 6,715; 6,885; 5,885 \\
Local split & 80\% train, 20\% holdout \\
Vocabulary / max length & 5,292 / 64 tokens \\
Encoder & $d=128$, 8 heads, 4 layers \\
Parameters & 1,499,163 \\
Dropout / optimizer & 0.15 / AdamW \\
Learning rate / weight decay & $3\times10^{-4}$ / $10^{-2}$ \\
Rounds / local epochs & 10 / 2 \\
Batch size & 16 \\
Nominal aggregation & FedTrimmedAvg, $\beta=0.1$ \\
Effective trim count & $b=0$ per tail \\
Malicious clients & 1 of 4 (client 3) \\
Poisoning & Gaussian, $\sigma=0.05$ \\
Seed / GPU & 42 / NVIDIA L40S \\
\bottomrule
\end{tabular}
\end{table}

\section{Results}

\subsection{Convergence}
The randomly initialized global model begins near chance accuracy (3.69\% versus $1/27=3.70\%$) with loss 3.3273. Accuracy rises to 56.93\% after one round and 74.87\% after five rounds. At round ten, federated-holdout accuracy is 77.20\% and loss is 0.5848 (Table~\ref{tab:convergence}). This trajectory demonstrates successful optimization despite one noisy upload in every round under the stated logical-federation setup.

\begin{table}[t]
\centering
\caption{Global-model convergence on the merged client holdout.}
\label{tab:convergence}
\begin{tabular}{rcc}
\toprule
Round & Cross-entropy loss & Accuracy (\%) \\
\midrule
0 & 3.3273 & 3.69 \\
1 & 1.3101 & 56.93 \\
2 & 0.8253 & 68.76 \\
3 & 0.7142 & 72.50 \\
5 & 0.6385 & 74.87 \\
7 & 0.6488 & 76.44 \\
9 & 0.6409 & 76.61 \\
10 & 0.5848 & 77.20 \\
\bottomrule
\end{tabular}
\end{table}

\subsection{Aggregate and Class-Resolved Detection}
At round ten, macro precision, recall, and F1 are 0.79, 0.78, and 0.77; weighted precision, recall, and F1 are 0.79, 0.77, and 0.77. The closeness of macro and weighted metrics is expected from the nearly balanced merged holdout. Table~\ref{tab:classes} reports all 27 class results.

Four attacks achieve 1.00 recall: false emergency vehicle, flooding DDoS, impossible kinematics, and Sybil. Benign recall is 0.99, indicating few benign messages are mislabeled in this holdout. Impersonation and pseudonym abuse reach F1 scores of 0.98 and 0.99. Conversely, heading manipulation (recall 0.40, F1 0.45) and position offset (recall 0.44, F1 0.57) are the weakest classes. Random position has recall 0.75 but precision 0.48, indicating substantial confusion into that class. Safety evaluation should prioritize these class-specific errors over aggregate accuracy.

\begin{table*}[t]
\centering
\caption{Round-10 class performance with one Gaussian-noise client. Recall is also the per-class accuracy reported by the pipeline.}
\label{tab:classes}
\resizebox{\textwidth}{!}{%
\begin{tabular}{lccc@{\hspace{1.2em}}lccc@{\hspace{1.2em}}lccc}
\toprule
Class & Precision & Recall & F1 & Class & Precision & Recall & F1 & Class & Precision & Recall & F1 \\
\midrule
Benign & 1.00 & 0.99 & 0.99 & Eventual stop & 0.59 & 0.53 & 0.56 & Impersonation & 0.99 & 0.98 & 0.98 \\
Constant position & 0.90 & 0.78 & 0.84 & False brake & 0.68 & 0.71 & 0.69 & Pseudonym abuse & 0.99 & 0.99 & 0.99 \\
Position offset & 0.80 & 0.44 & 0.57 & False emergency veh. & 1.00 & 1.00 & 1.00 & Flooding DDoS & 0.81 & 1.00 & 0.89 \\
Random position & 0.48 & 0.75 & 0.58 & False hazard & 0.54 & 0.72 & 0.62 & Beacon-rate abuse & 0.65 & 0.54 & 0.59 \\
Speed manipulation & 0.73 & 0.72 & 0.73 & Replay & 0.99 & 0.77 & 0.87 & GNSS spoofing & 0.65 & 0.77 & 0.70 \\
Acceleration manip. & 0.57 & 0.59 & 0.58 & Delayed message & 0.74 & 0.65 & 0.69 & Map-location spoof. & 0.87 & 0.77 & 0.82 \\
Heading manipulation & 0.52 & 0.40 & 0.45 & Timestamp shift & 0.88 & 0.94 & 0.91 & Ghost vehicle & 0.70 & 0.84 & 0.77 \\
Lane spoofing & 0.89 & 0.88 & 0.88 & Stale replay & 0.65 & 0.89 & 0.75 & False-object inject. & 0.97 & 0.70 & 0.81 \\
Impossible kinematics & 1.00 & 1.00 & 1.00 & Sybil & 1.00 & 1.00 & 1.00 & Object-position shift & 0.64 & 0.59 & 0.62 \\
\bottomrule
\end{tabular}}
\end{table*}

\subsection{Relation to the Honest Reference Run}
An earlier all-honest log reports 76.73\% accuracy, 0.77 macro F1, and loss 0.5825 on 5,402 holdout records. The noisy-client result is numerically similar. However, the earlier partition artifact was overwritten: its support and class counts differ from the preserved 5,400-record partition, while the repository documentation and current manifest record different Dirichlet concentrations. The runs are therefore not a paired control. The similarity is useful as a hypothesis---that this noise level may be tolerated---but the 0.47 percentage-point difference must not be interpreted as improvement or proof of robustness.

\section{Discussion}

\subsection{What the Experiment Establishes}
The experiment establishes three narrow facts. First, the implementation trains a nontrivial 27-class classifier: accuracy progresses from chance to 77.20\%. Second, training does not collapse when one client adds independent Gaussian noise with standard deviation 0.05 after each local update. Third, aggregate metrics hide substantial class variation, including weak heading and position-offset detection.

These observations are relevant to an IDS because the data contain benign and attack records regardless of the FL client-poisoning flag. The result is therefore not accuracy with no attacks; it is multi-attack classification under an additional attack on model training.

\subsection{What the Experiment Does Not Establish}
The experiment does not establish that temporal modeling improves detection: only tokens within one record are encoded. It does not establish trimmed-mean robustness because the effective trim count is zero. It does not establish superiority over the honest condition because the partition states differ. It does not establish generalization to unseen simulation runs, vehicles, densities, or attacks because evaluation randomly samples from the same client shards used for training. Finally, FL data locality alone does not imply differential privacy, secure aggregation, or resistance to inference attacks.

\subsection{Required Controlled Study}
A submission-grade evaluation should freeze one split manifest and reuse it for every method. At minimum, it should cross the following factors:
\begin{enumerate}
    \item aggregation: sample-weighted FedAvg, coordinate median, and trimmed mean with an integer trim count of at least one;
    \item adversary: none, Gaussian noise at multiple $\sigma$, sign flip, scaling, and an adaptive poisoning method \cite{fang2020local};
    \item malicious fraction: 0\%, 10\%, 20\%, and 25\%, with enough clients for each robust rule's assumptions;
    \item input: semantically selected observable fields, post-event proxies removed, longer non-truncating serialization, and numeric-feature baselines;
    \item context: single message versus fixed-length, sender-grouped temporal windows;
    \item evaluation: vehicle-, simulation-run-, and time-disjoint validation/test sets, with the canonical test set untouched until final selection.
\end{enumerate}
Each configuration should be repeated over at least five partition/training seeds. Report mean and 95\% confidence intervals for macro F1, balanced accuracy, benign false-positive rate, per-attack recall, expected calibration error, communication volume, and inference latency. Paired tests should compare methods on identical splits. A centralized upper bound, local-only models, and a non-transformer baseline are needed to isolate the contribution of federation and attention.

\section{Threats to Validity}
\textbf{Construct validity:} Serialization truncates most claimed-state fields, while early sender identifiers may enable simulation-specific memorization. Raw event-type values duplicate nine target names but are beyond the current cutoff. Some true-state fields may also be unavailable to a deployed detector. A feature-availability contract, explicit field selection, and sender/run-disjoint ablation are required.

\textbf{Internal validity:} The honest and noisy logs do not share a preserved partition fingerprint. Only one seed is available. The nominal robust strategy performs no trimming. These prevent causal attribution to the aggregator or attacker.

\textbf{External validity:} All clients run on one GPU host; no wireless, compute, or dropout heterogeneity is represented. The corpus is simulator-derived and balanced, unlike operational traffic. No unseen attack, unseen map, or cross-dataset evaluation is reported.

\textbf{Statistical conclusion validity:} The current paper reports descriptive metrics from one run. No confidence interval or significance test is meaningful until repeated matched trials are available.

\textbf{Reproducibility:} The pipeline, manifests, random seed, software versions, and logs are preserved. Before archival release, generated payload artifacts should be separated from learning data, environment dependencies pinned, and each result linked to immutable hashes of the source, split, vocabulary, and configuration.

\section{Conclusion}
This paper audited a four-client transformer IDS over 27 VANET-IDS26 classes under one Gaussian-noise client. The global model reached 77.20\% accuracy and 0.77 macro F1 on a federated holdout, with strong detection of several conspicuous attacks but weak recall for heading manipulation and position offset. The audit changes the scientific interpretation: the configured trim proportion removes no client update, the model has no cross-message temporal context, the token cutoff removes most claimed-state and all event fields, and the honest reference run is not a matched control. Raw event names do not reach the current encoder, although they remain a latent leakage risk if the input length changes. These findings turn a promising run into a precise experimental baseline. A controlled, feature-audited, group-disjoint, multi-seed study with an effective robust aggregator is the necessary next step toward a defensible claim of robust temporal federated intrusion detection.

\section*{Data and Code Availability}
The implementation, manifests, and run logs are maintained in the accompanying VANET-IDS26 repository. The evaluated dataset is stored in the following canonical local files:
\begin{itemize}
    \item \path{/opt/sahsan03/VANET-IDS26/data/vanet_ids26_master.csv};
    \item \path{/opt/sahsan03/VANET-IDS26/data/vanet_ids26_train.csv};
    \item \path{/opt/sahsan03/VANET-IDS26/data/vanet_ids26_validation.csv}; and
    \item \path{/opt/sahsan03/VANET-IDS26/data/vanet_ids26_test.csv}.
\end{itemize}
The public release will include an immutable commit identifier, a dataset access statement, and scripts that regenerate every table. \revisionnote{Replace the local paths with the archival dataset DOI and license before submission.}

\section*{Acknowledgment}
\revisionnote{Insert funding, compute-resource, dataset-creator, and contributor acknowledgments after de-anonymization.}

\bibliographystyle{IEEEtran}
\bibliography{references}

\end{document}
